\chapter{Resonances and Multiple Channels}

\marginnoteinline{Project goal at the top of the first page: develop a numerical scheme for evaluating optical models, including volume, surface, and spin--orbit terms.}

\sourcepage{1}
\section{The collision matrix and the logarithmic derivative}
From the preceding discussion,
\begin{equation}
  S_\ell
  =\frac{1-R_\ell\widehat S_\ell+\ii R_\ell\widehat P_\ell}
         {1-R_\ell\widehat S_\ell-\ii R_\ell\widehat P_\ell}
   \ee^{2\ii\Delta_\ell}
  =\ee^{2\ii(\alpha_\ell+\Delta_\ell)}.
  \label{eq:res-S-R}
\end{equation}
Here $\widehat S_\ell$ is the shift factor, $\widehat P_\ell$ is the penetrability, and $\Delta_\ell$ is the hard-sphere phase.  The partial elastic and reaction cross sections are
\begin{align}
  \sigma_{\mathrm{el},\ell}
  &=\pi\bar\lambda^2(2\ell+1)|1-S_\ell|^2,
  \label{eq:res-elastic-partial}\\
  \sigma_{r,\ell}
  &=\pi\bar\lambda^2(2\ell+1)(1-|S_\ell|^2).
  \label{eq:res-reaction-partial}
\end{align}

Define the logarithmic derivative
\[
  D_\ell\equiv\frac{1}{R_\ell}.
\]
Then
\begin{equation}
  S_\ell
  =\frac{D_\ell-\widehat S_\ell+\ii\widehat P_\ell}
         {D_\ell-\widehat S_\ell-\ii\widehat P_\ell}
   \ee^{2\ii\Delta_\ell}.
  \label{eq:res-S-D}
\end{equation}
Consequently,
\begin{equation}
  1-S_\ell
  =\ee^{2\ii\Delta_\ell}
  \left[
    \left(\ee^{-2\ii\Delta_\ell}-1\right)
    -\frac{2\ii\widehat P_\ell}
    {D_\ell-\widehat S_\ell-\ii\widehat P_\ell}
  \right].
  \label{eq:res-amplitude-split}
\end{equation}
The two terms exhibit the possible interference between the potential-scattering amplitude and the compound or direct-reaction amplitude.

Writing
\[
  D_\ell=D_\ell^{\mathrm{Re}}+\ii D_\ell^{\mathrm{Im}},
\]
the modulus of the collision matrix follows from
\begin{equation}
  |S_\ell|^2
  =\frac{(D_\ell^{\mathrm{Re}}-\widehat S_\ell)^2
          +(D_\ell^{\mathrm{Im}}+\widehat P_\ell)^2}
         {(D_\ell^{\mathrm{Re}}-\widehat S_\ell)^2
          +(D_\ell^{\mathrm{Im}}-\widehat P_\ell)^2}.
  \label{eq:res-S-modulus}
\end{equation}

\sourcepage{2}
Therefore
\begin{equation}
  \boxed{
  1-|S_\ell|^2
  =\frac{-4D_\ell^{\mathrm{Im}}\widehat P_\ell}
  {(D_\ell^{\mathrm{Re}}-\widehat S_\ell)^2
   +(D_\ell^{\mathrm{Im}}-\widehat P_\ell)^2}.}
  \label{eq:res-reaction-factor}
\end{equation}
Let the formal resonance energy $E_0'$ be defined by
\begin{equation}
  \operatorname{Re}D_\ell(E_0')=0.
  \label{eq:res-formal-energy}
\end{equation}
Near this energy,
\begin{equation}
  \operatorname{Re}D_\ell(E)
  \simeq
  (E-E_0')
  \left.
    \frac{\dd}{\dd E}\operatorname{Re}D_\ell(E)
  \right|_{E=E_0'}.
  \label{eq:res-D-linearization}
\end{equation}

Why does this condition identify a resonance?  For $\ell=0$,
\[
  \widehat S_0=0,
  \qquad
  \widehat P_0=kr.
\]
Both the elastic expression based on Eq.~\eqref{eq:res-amplitude-split} and the reaction factor in Eq.~\eqref{eq:res-reaction-factor} are maximized when $D_0^{\mathrm{Re}}=0$.  The page emphasizes that this simple conclusion is not generally true for $\ell>0$, where the shift and penetrability have additional energy dependence.

\sourcepage{3}
\section{Wave-function picture of a resonance}
The source sketches several interior and exterior radial wave functions separated by the channel radius $r=a$.  At a resonance, an interior eigenfunction happens to join smoothly to an exterior solution near an extremum, so that the logarithmic derivative at the boundary approaches zero:
\[
  \left.\frac{\psi'}{\psi}\right|_{r=a}\longrightarrow0.
\]
The drawing contrasts nonresonant matches with the resonant case in which the interior oscillation has the appropriate phase at the boundary.

\begin{center}
\begin{tikzpicture}[x=1cm,y=1cm,>=Latex]
  \draw[->] (0,0) -- (10,0) node[right] {$r$};
  \draw[dashed] (4.5,-2.0) -- (4.5,2.1) node[above] {$r=a$};
  \node at (2.2,1.8) {interior};
  \node at (7.0,1.8) {exterior};
  \draw[domain=0.2:4.5,samples=140,smooth]
    plot (\x,{0.72*sin(155*\x)});
  \draw[domain=4.5:9.2,samples=140,smooth]
    plot (\x,{0.8*cos(42*(\x-4.5)) * exp(-0.13*(\x-4.5))});
  \fill (4.5,{0.72*sin(155*4.5)}) circle (1.6pt);
  \draw[->] (7.5,-1.25) -- (4.75,-0.55);
  \node[align=left] at (8.0,-1.45) {$\psi'/\psi\to0$\\at the match};
\end{tikzpicture}
\end{center}
\editorial{The TikZ figure is a schematic transcription of the three hand-drawn wave-function examples rather than a digitization of their exact curves.}

\sourcepage{4}
\section{Neutron width, reaction width, and observed energy}
Define
\begin{align}
  \Gamma_n
  &\equiv
  \frac{-2\widehat P_\ell}
  {\left.
    \dfrac{\dd}{\dd E}\operatorname{Re}D_\ell(E)
   \right|_{E=E_0'}},
  \label{eq:res-neutron-width}\\
  \Gamma_r
  &\equiv
  \frac{2\operatorname{Im}D_\ell}
  {\left.
    \dfrac{\dd}{\dd E}\operatorname{Re}D_\ell(E)
   \right|_{E=E_0'}},
  \label{eq:res-reaction-width}\\
  \Gamma&\equiv\Gamma_n+\Gamma_r.
  \label{eq:res-total-width}
\end{align}
The observed resonance energy is shifted from the formal resonance energy:
\begin{equation}
  E_0
  =E_0'
  +\frac{\widehat S_\ell}
  {\left.
    \dfrac{\dd}{\dd E}\operatorname{Re}D_\ell(E)
   \right|_{E=E_0'}}.
  \label{eq:res-observed-energy}
\end{equation}
Thus the formal resonance energy is inferred from the observed resonance after accounting for the shift factor.  The page also notes that $\sigma_r$ is proportional to the penetrability $\widehat P_\ell$.

With these definitions, the elastic cross section is written
\begin{equation}
  \sigma_{\mathrm{el},\ell}
  =\pi\bar\lambda^2(2\ell+1)
  \left|
    \ee^{2\ii\Delta_\ell}
    \left[
      \left(\ee^{-2\ii\Delta_\ell}-1\right)
      +\frac{\ii\Gamma_n}
       {E-E_0+\ii\Gamma/2}
    \right]
  \right|^2,
  \label{eq:res-elastic-bw}
\end{equation}
and the reaction cross section becomes
\begin{equation}
  \sigma_{r,\ell}
  =\pi\bar\lambda^2(2\ell+1)
  \frac{\Gamma_n\Gamma_r}
  {(E-E_0)^2+(\Gamma/2)^2}.
  \label{eq:res-reaction-bw}
\end{equation}
If the compound-nucleus formation cross section is
\[
  \sigma_{c,\ell}
  =\pi\bar\lambda^2(2\ell+1)
  \frac{\Gamma_n\Gamma}
  {(E-E_0)^2+(\Gamma/2)^2},
\]
then
\begin{equation}
  \sigma_{r,\ell}=\sigma_{c,\ell}\frac{\Gamma_r}{\Gamma}.
  \label{eq:res-branching-xsec}
\end{equation}

\sourcepage{5}
\section{Meaning of the resonance width}
The widths $\Gamma$ have units of energy.  Earlier they were connected to the lifetime before decay; the page asks what else $\Gamma$ means and records
\[
  \Gamma\propto\frac{1}{\tau_c}.
\]

Consider the Lorentzian factor in Eq.~\eqref{eq:res-reaction-bw}.  Its maximum occurs at $E=E_0$.  When
\[
  (E-E_0)^2=(\Gamma/2)^2,
\]
the cross section is one-half of its maximum value.  Hence
\[
  |E-E_0|=\frac{\Gamma}{2},
\]
and the full width at half maximum is
\begin{equation}
  \boxed{\mathrm{FWHM}=\Gamma.}
  \label{eq:res-fwhm}
\end{equation}

The reduced width is defined by
\begin{equation}
  \boxed{
  \gamma^2
  \equiv
  -\left[
  \left.
    \frac{\dd}{\dd E}\operatorname{Re}D_\ell(E)
  \right|_{E=E_0'}
  \right]^{-1}
  =\frac{\Gamma_n}{2\widehat P_\ell}.}
  \label{eq:res-reduced-width}
\end{equation}
The handwritten comment calls this the ``reduced width'': the neutron width reduced by the penetrability factor.

\sourcepage{6}
\section{Multichannel scattering and angular-momentum coupling}
Consider compound-nucleus formation,
\[
  a+X\longrightarrow C,
\]
with relative orbital angular momentum $L$, projectile spin $I_p$, and target spin $I_t$.  First couple the intrinsic spins to form the channel spin,
\begin{equation}
  \bm I_p+\bm I_t=\bm S,
  \label{eq:res-channel-spin}
\end{equation}
and then couple channel spin and orbital angular momentum to total angular momentum,
\begin{equation}
  \bm S+\bm L=\bm J.
  \label{eq:res-total-J-coupling}
\end{equation}
Here $S$ is the channel spin and $J$ is the total angular momentum.

An alternative coupling scheme first forms a projectile orbital angular momentum,
\[
  \bm I_p+\bm L=\bm J_p,
  \qquad
  \bm J_p+\bm I_t=\bm J.
\]
The conservation statements written on the page are
\begin{align*}
  \text{total angular momentum:}\quad
  &\bm J=\bm I_p+\bm I_t+\bm L
   =\bm I_p'+\bm I_t'+\bm L',\\
  \text{parity:}\quad
  &\pi=\pi_p\pi_t(-1)^L
  =\pi_p'\pi_t'(-1)^{L'},\\
  \text{azimuthal quantum number:}\quad
  &m_J=-J,-J+1,\ldots,J.
\end{align*}

\sourcepage{7}
\section{Example: neutron entrance channels}
Take a neutron incident on a target with
\[
  I_t^\pi=\left(\frac32\right)^-
\]
to form a compound state with
\[
  J^\pi=2^+.
\]
The neutron has spin $I_p=1/2$, so the possible channel spins are
\[
  S=1,2.
\]
For a given $S$, the triangle rule is
\begin{equation}
  |J-S|\le L\le J+S.
  \label{eq:res-L-triangle}
\end{equation}
Thus
\begin{align*}
  S=1:&\qquad 1\le L\le3,\\
  S=2:&\qquad 0\le L\le4.
\end{align*}
The source page uses check marks and crosses to test candidate $L$ values against these bounds.

Parity conservation gives
\[
  \pi=\pi_p\pi_t(-1)^L
  =(+1)(-1)(-1)^L
  =(-1)^{L+1}.
\]
Since the compound nucleus has positive parity,
\begin{equation}
  L=1,3
  \label{eq:res-example-odd-L}
\end{equation}
for the candidate range shown.  For spinless projectile and target, $S=0$, so $J=L$; the page remarks that all $L$ values are then in principle applicable.

\editorial{A few of the handwritten row labels/check marks in the channel table are visually ambiguous.  The inequalities and parity conclusion above are transcribed from the clearly legible statements surrounding the table.}

\sourcepage{8}
\section{Statistical spin factor}
Given $I_p$, $I_t$, and $L$, there are
\[
  (2I_p+1)(2I_t+1)(2L+1)
\]
uncoupled magnetic substates.  Only $2J+1$ of the coupled states correspond to a specified $J$.  Define the statistical spin factor
\begin{equation}
  \boxed{
  g_J=\frac{2J+1}{(2I_p+1)(2I_t+1)}.}
  \label{eq:res-statistical-spin-factor}
\end{equation}
The factor $(2L+1)$ is omitted from $g_J$ because it replaces the partial-wave degeneracy already appearing in expressions such as $(2\ell+1)\bar\lambda^2$.

For spinless projectile and target, the preceding single-resonance formation expression was
\[
  \sigma_{c,\ell}
  =\pi\bar\lambda^2(2\ell+1)
  \frac{\Gamma_n\Gamma}
  {(E-E_0)^2+(\Gamma/2)^2}.
\]
When conservation of $J$ and the particle spins are included, replace it by
\begin{equation}
  \boxed{
  \sigma_{c,J}
  =\pi\bar\lambda^2 g_J
  \frac{\Gamma_n\Gamma}
  {(E-E_0)^2+(\Gamma/2)^2}.}
  \label{eq:res-spin-dependent-formation}
\end{equation}
Thus $g_J$ depends on $J$ and on the projectile and target spins.  The page closes with the prompt: think about the effect of conservation on the decay scheme.

\sourcepage{9}
\section{Channel functions and the multichannel \texorpdfstring{$S$}{S} matrix}
For a single spinless channel, the angular and radial ingredients are
\[
  P_\ell(\cos\theta)
  \qquad\text{and}\qquad
  h_\ell^\pm(kr).
\]
For a multichannel problem, label a channel by
\begin{equation}
  c=[\alpha,L,S,J,m_J],
  \label{eq:res-channel-label}
\end{equation}
where $\alpha$ denotes the particle arrangement and any additional quantum numbers.  The channel wave function contains a spin-angular function
\[
  \mathcal{Y}_{LS}^{Jm_J}
\]
built by coupling a spherical harmonic $Y_L^{m_L}(\theta,\phi)$ to a spin function $\chi_S^{m_S}$, together with the radial function $h_L^\pm$.

Explicitly, the coupling indicated by the diagram may be written
\begin{equation}
  \mathcal{Y}_{LS}^{Jm_J}(\widehat{\bm r})
  =\sum_{m_L,m_S}
  \langle Lm_L,Sm_S\mid Jm_J\rangle
  Y_L^{m_L}(\widehat{\bm r})\chi_S^{m_S}.
  \label{eq:res-spin-angular-function}
\end{equation}
\editorial{Equation~\eqref{eq:res-spin-angular-function} makes explicit the Clebsch--Gordan coupling shown pictorially in the scan.}
The result is a multichannel collision matrix $S_{cc'}$.

As a simple $m=0$ illustration with $L=0$ and $L=1$, the notes write coupled radial equations of the form
\begin{align}
  \left(-\frac{\hbar^2}{2m}\frac{\dd^2}{\dd r^2}-E\right)\psi_0
  +V_{00}\psi_0+V_{01}\psi_1&=0,
  \label{eq:res-coupled-l0}\\
  \left(-\frac{\hbar^2}{2m}\frac{\dd^2}{\dd r^2}
  +\frac{\hbar^2L(L+1)}{2mr^2}-E\right)\psi_1
  +V_{11}\psi_1+V_{10}\psi_0&=0,
  \qquad (L=1).
  \label{eq:res-coupled-l1}
\end{align}
The angular coupling matrix is defined on the page by
\begin{equation}
  V_{LL'}
  =\int_{-1}^{1}P_L(\mu)V(\mu)P_{L'}(\mu)\,\dd\mu.
  \label{eq:res-angular-potential-matrix}
\end{equation}
